\section{The characterwise estimate on the minus side}
\label{sec:minus-characterwise}

Let \(E\subseteq\mathbb Q(\mu_N)\) be a CM field. Put
\(D=[E:\mathbb Q]\), \(G=\operatorname{Gal}(E/\mathbb Q)\), and
let \(\ell\nmid 2D\) be prime. Let \(\mathcal O\) be the ring of
integers of a finite extension of \(\mathbb Q_\ell\) in the chosen
algebraic closure, with uniformizer \(\varpi\), valuation
\(v(\varpi)=1\), and \(r_{\mathcal O}=v(\ell)\).
Write \(W=\operatorname{Frac}(\mathcal O)/\mathcal O\). A character
\(\psi:G\to\mathcal O^\times\) pulls back to a character \(\chi\) of
level \(N\). Its order divides \(D\) and is prime to \(\ell\).
Fix \(e\ge1\), put \(j=e+1\), and assume \(\chi(-1)=(-1)^j\).
Write \(\kappa\) for the \(\ell\)-adic cyclotomic character and put
\(\alpha=(\chi\kappa^e)^{-1}\), using the Galois character attached
to \(\chi\). Let \(V(\alpha)\) be the \(\mathcal O\)-length of the
strict subgroup of \(H^1(\mathbb Q,W(\alpha))\), with zero restriction
to all inertia groups and to the decomposition groups at \(\ell\).
The source proves
\begin{equation}
    V(\alpha)\le v\bigl(\alpha_N(j,\chi)\bigr).
    \label{eq:characterwise-bern}
\end{equation}
On the right, \(\alpha_N\) is the integral generalized Bernoulli
number defined in Section~\ref{sec:odd-prime-minus}. We give the
conductor and local conditions, the finite descent, its exceptional
case, and the passage to a bound on the \(E\)-group. We specify below
the cyclotomic derivative-system construction used by the descent.

\subsection{Conductor and local conditions}

Write \(f=\operatorname{cond}\chi\) and \(\xi=\chi^*\). Raising the
level from \(f\) to \(N\) preserves the Galois character, its order
and its parity. By the primitive-value convention,
\[
  \alpha_N(j,\chi)=(N/f)\alpha_f(j,\xi).
\]
The factor \(N/f\) is integral, so it suffices to prove
\eqref{eq:characterwise-bern} at level \(f\).

If \(p\mid f\), \(p\ne\ell\), inertia at \(p\) contains an element
on which \(\xi\) is nontrivial. Indeed the inertia image in the units
modulo \(f\) contains the full \(p\)-component, and triviality on
this component would lower the conductor. A nontrivial root of unity
of order prime to \(\ell\) differs from \(1\) by a unit in
\(\mathcal O\), since reduction is injective on these roots. The
cyclotomic character is unramified away from \(\ell\). Thus at \(p\)
some \(\tau\in I_p\) has \(\alpha(\tau)-1\) a unit. An
inertia coboundary is then a decomposition-group coboundary.
Indeed subtract a local coboundary so the cocycle \(c\) vanishes on
inertia. For \(g\) in the decomposition group, compare \(c(\tau g)\)
with \(c(g(g^{-1}\tau g))\). This gives
\((\alpha(\tau)-1)c(g)=0\), hence \(c(g)=0\).

Choose a lift \(\delta\) of complex conjugation. The parity gives
\(\alpha(\delta)=-1\), so \(\alpha(\delta)-1\) is a unit. Each
class has a unique representative \(c\) satisfying \(c(\delta)=0\):
subtract the coboundary of
\((\alpha(\delta)-1)^{-1}c(\delta)\).
For \(a_0\ne0\), denote by \(S(a_0)\) the module of these normalized
continuous cocycles killed pointwise by \(a_0\), vanishing on inertia
outside \(f\ell\), and coboundaries at the places dividing \(f\ell\)
with local primitives killed by \(a_0\). Every strict class belongs
to the image of one of these modules. The finite image of \(c\) and
one chosen local primitive above \(\ell\) have a common annihilator.
At conjugate places above \(\ell\) the new primitive has the form
\(\alpha(g)m-c(g)\) and is killed by the same annihilator. At a
ramified prime \(p\ne\ell\) the equation
\((\alpha(\tau)-1)m=c(\tau)\) gives the same property by the unit
just found. Outside \(f\ell\), inertia acts trivially, so a local
coboundary is zero. The modules with \(a_0=\varpi^h\) increase with
\(h\). A uniform finite upper bound for their lengths bounds the union:
a longer finite strict chain in the union would have a finite set of
witnesses already contained in one term.

\subsection{The residual cyclotomic exception}

Let \(\overline{\alpha}\) be reduction modulo \(\varpi\). If
\(\overline\alpha=\overline\kappa\), the preceding inertia observation
forces \(\alpha\) to be unramified outside \(\ell\): away from \(\ell\)
its inertia values have order prime to \(\ell\) and reduce to one.
Take \(t\) in inertia at a place over \(\ell\) which acts by inversion
on \(\mu_\ell\). Such \(t\) exists by total ramification of the local
cyclotomic extension. Normalize at \(t\) in this case.
Since \(\alpha(t)-1\) is a unit, vanishing of the cocycle at \(t\)
makes its local primitive at the chosen place zero.

Suppose a normalized cocycle is killed by \(\varpi\). Every
\(\mathbb F_\ell\)-linear coordinate on \(W[\varpi]\) takes it to a
cocycle with values in \(\mathbb F_\ell(1)\): the residual action is
multiplication by \(\overline\kappa\), which lies in the prime field.
Identify this target with \(\mu_\ell\). Kummer theory writes the
coordinate cocycle as \(\sigma\mapsto\sigma\beta/\beta\) with
\(a=\beta^\ell\in\mathbb Q^\times\). Its vanishing on inertia at
a place above each rational prime shows \(\ell\mid v_p(a)\):
\(\beta\) is in the local unramified extension there, whose value
group is \(\mathbb Z\). The sign of \(a\) is an \(\ell\)-th power
as \(\ell\) is odd. Rational prime factorization gives
\(a=b^\ell\) for \(b\in\mathbb Q^\times\).
Now \(\beta/b\in\mu_\ell\). Normalization at \(t\), which acts by
inversion, forces \(\beta/b=1\). Every coordinate of the cocycle
therefore vanishes. They separate points, so the cocycle vanishes.
Applying this argument to successive powers of \(\varpi\) proves
\(S(a_0)=0\) for all \(a_0\). This proves
\eqref{eq:characterwise-bern} in the exceptional branch.

\subsection{The normalized finite Stickelberger value}

Assume \(\overline\alpha\ne\overline\kappa\). For a fixed \(a_0\ne0\)
take \(m\ge\max(1,v_\ell(f))\), and set
\(P=f\ell^{m-v_\ell(f)}\). The sets of primes dividing \(P\) and
\(f\ell\) agree. Choose a lift \(\sigma_t\) of each
\(t\in(\mathbb Z/P\mathbb Z)^\times\). Write \([t]\) for the least
nonnegative representative and choose \(b\) a unit representative
modulo \(P\) coming from a \(g\) with
\(\overline\alpha(g)\ne\overline\kappa(g)\). Then
\(\xi(b)b^{e+1}-1\) is a unit, since multiplication of
\(\alpha(g)-\kappa(g)\) by \(-\alpha(g)^{-1}\) gives the same
quantity modulo \(\varpi\). Define
\[
 h_b(t)=\left\lfloor\frac{b[t]}P\right\rfloor-
        \left\lfloor\frac{b[-t]}P\right\rfloor,\qquad
 z=\sum_t h_b(t)\alpha(\sigma_t)^{-1}.
\]
The level-\(P\) values of \(\alpha(\sigma_t)^{-1}\) are congruent
modulo \(\ell^m\) to \([t]^e\xi([t])\). Write \(S\) for the sum
obtained by this replacement. Then \(|z-S|_\ell\le|\ell^m|_\ell\).

Here is the source's normalization of the other congruence. Put
\(d_j=2j\prod_{p\le j+1}p\) over primes and
\(B=B_{j,\xi_P}\), the generalized Bernoulli number of the
level-\(P\) character induced by \(\xi\). Then
\[
 \left|d_j\left(b^e\xi(b)S
        -2\{\xi(b)b^j-1\}\frac{B}{j}\right)\right|_\ell
        \le|\ell^m|_\ell .
\]
For completeness take \(u=[t]\), \(h=\lfloor bu/P\rfloor\),
\(y=bu-hP\). Both \(u,y\) are units modulo \(P\) and
\(h_b(t)=2h+1-b\). The polynomial
\(C(x)=P^{j-1}B_j(x/P)/j\) satisfies
\[
       C(x)=\frac{x^j}{jP}-\frac{x^{j-1}}2+
                        \frac{2P}{d_j}Z_x,\qquad Z_x\in\mathbb Z.
\]
This follows by the denominator calculation of
Section~\ref{sec:odd-prime-minus} applied to the terms of Bernoulli
index at least two. Expand \(y^j=(bu-hP)^j\) through its linear
term in \(P\), and \(y^{j-1}\) through its constant term.
The leading terms then cancel in
\[
 d_j\{b^e h_b(t)u^e+2(C(y)-b^j C(u))\}\in P\mathbb Z.
\]
Indeed their coefficient of \(b^e u^e\) is
\(h_b(t)-2h-1+b=0\), and \(j\mid d_j\) clears the remaining
binomial denominators. Multiply this relation by \(\xi_P(bt)\),
sum over units, and permute them by \(t\mapsto bt\). This gives
the asserted congruence.

Raising the level from \(f\) to \(P\) adds at most the Euler factor
\(1-\xi(\ell)\ell^e\), of \(\ell\)-adic norm one because \(e\ge1\).
Factors at the primes already dividing \(f\) are one. Consequently
\(|B/j|_\ell=|B_{j,\xi}/j|_\ell>0\). Choose \(m\) large enough that
\[
 |\ell^m|_\ell
     < |a_0|_\ell\,|d_j|_\ell\,|B_{j,\xi}/j|_\ell .
\]
This is possible before choosing \(b\) since the right side is fixed.
The elements \(b^e\xi(b)\), \(2\), and \(\xi(b)b^j-1\) are units.
The strong triangle inequality in the two congruences gives
\(|z|_\ell=|B_{j,\xi}/j|_\ell\). It follows that \(z\ne0\),
\(a_0z\mid\ell^m\), and \(z\mid\alpha_f(j,\xi)\) in \(\mathcal O\):
the last ratio has norm at most one because
\(\alpha_f(j,\xi)=f(\prod_{p\le j+1}p)B_{j,\xi}\).

\subsection{Local reciprocity and the finite descent}

The reciprocity used in the source is a consequence of the
finite-module Poitou--Tate sequence and the unramified-subgroup
theorem in local Tate duality. In the conventions of Milne's
\emph{Arithmetic Duality Theorems}, second edition, these are
I.4.10 and I.2.6 respectively. The local assertion applies at an
auxiliary prime \(q\ne\ell\) to unramified finite coefficients of
\(\ell\)-power order. The global assertion uses the sum of the
local invariant maps on \(H^2(\mu_{\ell^h})\). The archimedean
contribution vanishes for odd \(\ell\).\footnote{\href{https://www.jmilne.org/math/Books/ADTnot.pdf}{Milne, second edition, Chapter I, Theorems 2.6 and 4.10}.}

More explicitly, take a finite coefficient quotient \(R=\mathcal
O/(a)\), killed by \(\ell^h\), and cocycles \(c,n\) with characters
\(\alpha,\kappa\alpha^{-1}\) modulo \(a\). Assume their inertia
restrictions vanish outside \(N'q\), where \(\ell\mid N'\), the
first cocycle is a coboundary at the places dividing \(N'\) away
from \(q\), both characters are trivial modulo \(a\) at \(q\), and
\(c\) vanishes on \(I_q\). For each \(\mathbb Z/\ell^h\)-linear
functional \(\lambda:R\to\mathbb Z/\ell^h\), the cyclotomic
two-cocycle
\[
        (g,h')\longmapsto\lambda(n(g)\alpha(g)c(h'))
\]
has zero local class away from \(q\). At the primes outside the
ramification set it is an unramified cup product. At the primes
in \(N'\) away from \(q\) one factor is a coboundary. The global
invariant sum and injectivity of the local invariant map therefore
give zero local class at \(q\). Let \(\phi\) be a Frobenius at \(q\).
For an unramified \(\mathbb Z/\ell^h\)-valued homomorphism \(x\),
one has \(x(h')c(\phi)=x(\phi)c(h')\) on the decomposition group:
both sides are additive, coincide at \(\phi\), and vanish on
inertia. Consequently the cup product of
\(\lambda(n(\cdot)c(\phi))\) with \(x\) is the preceding zero local
class multiplied by \(x(\phi)\). Unramified orthogonality puts
\(\lambda(n(\cdot)c(\phi))\) in the unramified subgroup. Varying
\(\lambda\) now gives \(n(\tau)c(\phi)=0\) for all \(\tau\in I_q\).
These functionals separate points by decomposing the finite
\(\ell\)-group into cyclic summands. This is the local
annihilation relation used in the descent.

Choose an auxiliary prime \(r\equiv1\bmod P\), and let \(L\) range
over lists of distinct primes \(q\ne r\) with \(q\equiv1\bmod P\)
and \(\ell\nmid\operatorname{ord}_{q}(r)\).
Write \(H=G_{\mathbb Q(\mu_P)}\). The source constructs
\(\nu_L|_H:H\to W[\ell^m]\), continuous additive and equivariant
under conjugation by the dual character \(\kappa\alpha^{-1}\).
Its inertia vanishes away from \(Pr\prod_{q\in L}q\).
For a fixed generator \(w\) of \(W[\ell^m]\), the initial value on a
suitable inertia element at \(r\) is \(zw\) times a unit. At each
admissible new prime \(q\), a suitable inertia value of \(\nu_{qL}\)
is the Frobenius value of \(\nu_L\) times a unit. Here \(qL\) means
the list with \(q\) added. The construction starts with cyclotomic
elements with prescribed Stickelberger ideal factorizations, applies
the multiplicative derivatives
\(\prod_{i=0}^{q-2}(\tau_q^i x)^i\) for generators \(\tau_q\) of the
new cyclotomic factors, and takes Kummer classes. The relations on
these derivatives give the stated initial and transition values.

Write \((A_L)=\operatorname{Ann}_{\mathcal O}(\nu_L(H))\) and
\(J=v(\ell^m)\). The coefficient ideal contains \(\ell^m\), so
\(j_L=v(A_L)\le J\). For \(a_0\mid A_L\) let \(S_L\subset S(a_0)\)
consist of normalized cocycles vanishing on the decomposition groups
at the primes dividing \(r\prod_{q\in L}q\). For \(S_L\ne0\),
choose \(c\ne0\) in it. The source's Chebotarev selection supplies
an admissible \(q\) outside this finite list and a Frobenius
\(\phi\in H\) such that \(c(\phi)\ne0\) and \(\nu_L(\phi)\)
generates the \(\mathcal O\)-span of \(\nu_L(H)\). Thus
the new inertia value has annihilator \((A_L)\).
Since inertia at \(q\) lies in \(H\), we have
\(A_{qL}=A_L t\) for some nonzero \(t\in\mathcal O\), up to a unit.

The equivariant derivative has a global extension with controlled
image. Write \(\beta=\kappa\alpha^{-1}\). The assumption
\(\overline\alpha\ne\overline\kappa\) supplies \(d\in G_{\mathbb Q}\)
with \(\beta(d)-1\) a unit. Given any of the equivariant derivative
homomorphisms \(\nu\) on \(H\), define
\[
 n(g)=-(\beta(d)-1)^{-1}\nu([g,d]),\qquad
 [g,d]=gdg^{-1}d^{-1}.
\]
Every such commutator belongs to \(H\), since \(G_{\mathbb Q}/H\)
is abelian. This function is locally constant by continuity of the
commutator and local constancy of \(\nu\) on \(H\). The identity
\([gh,d]=g[h,d]g^{-1}[g,d]\) and equivariance give
\(n(gh)=n(g)+\beta(g)n(h)\). For \(u\in H\),
\(\nu([u,d])=(1-\beta(d))\nu(u)\), whence \(n(u)=\nu(u)\).
All values of \(n\) are multiples of values in \(\nu(H)\).
In particular the annihilator of \(\nu(H)\) kills the whole cocycle.
Away from \(P\), inertia lies in \(H\), so extension preserves the
required inertia conditions. The same observation identifies the
inertia values in the local comparison at the selected \(q\). This
proves the extension used with the duality relation above.
\footnote{\href{https://github.com/Deicyde/riemann-hypothesis/blob/40e56bab026552e089d7afcfeff08d2defb9c005/src/Development/AutomorphicArithmetic/AutomorphicVarieties/PrimitiveFrobeniusFixed/ImportsAlt.lean}{The commutator formula and its restriction and image properties}.}

To see the resulting cost, choose a generator of \(W[A_{qL}]\)
and write the new inertia value as \(y\) times this generator. If
the value is nonzero, \(v(y)=v(t)\) by its annihilator; if it is zero,
then \(A_L\) is a unit, and \(a_0\mid A_L\) forces \(S_L=0\).
Using the extension of the equivariant derivative to a dual cocycle
in the source's reciprocity step gives \(y\,c'(\phi)=0\) for every
\(c'\in S_L\). Thus \(t\) kills the image of evaluation on \(S_L\)
at \(\phi\), which has length at most \(v(t)\) in \(W\). Its kernel
is contained in \(S_{qL}\): the cocycles vanish on inertia at \(q\)
and the unramified quotient is topologically generated by Frobenius.
Conjugacy propagates the vanishing to the other places above \(q\);
the coefficient character acts trivially modulo \(a_0\) throughout
these decomposition groups. We obtain
\[
   \operatorname{length}S_L+j_L
       \le \operatorname{length}S_{qL}+j_{qL}.
\]
Since some \(c(\phi)\ne0\), \(t\) is a nonunit and \(j_{qL}>j_L\).
Since \(j_L\le J\), iteration ends with a zero submodule. Hence
\(\operatorname{length}S_L+j_L\le J\) on every such list.

The first step at \(r\) begins with the whole \(S(a_0)\).
Put \(k_0=\ell^m/z\). The initial inertia value \(zw\) has
annihilator \((k_0)\), and \(a_0\mid k_0\). The same evaluation gives
\(A_{\varnothing}=k_0t_0\) up to a unit and
\(\operatorname{length}S(a_0)\le
  \operatorname{length}S_{\varnothing}+v(t_0)\).
Combining this with the bound for the empty list gives
\[
  \operatorname{length}S(a_0)
        \le J-v(k_0)=v(z)\le v(\alpha_f(j,\xi)).
\]
This bound is independent of \(a_0\). Taking the increasing union
described above gives \eqref{eq:characterwise-bern}.

\subsection{Return to the minus eigenspace}

Take \(\mathcal O\) large enough to contain all character values
of \(G\). Since \(\ell\nmid D\), restriction to \(G_E\) identifies
the twisted strict group over \(\mathbb Q\) with the
\(\psi\)-eigenspace over \(E\). Indeed positive cohomology of \(G\)
on the \(\ell\)-primary invariants vanishes by averaging over \(G\);
the same restriction and averaging applies to the local subgroups,
whose indices divide \(D\). Untwisting by \(\psi\) produces the
indicated eigenspace. In the sign of complex conjugation under
consideration, the characters are those with
\(\psi(c)=(-1)^{e+1}\). Base change multiplies a finite
\(\mathbb Z_\ell\)-length by \(r_{\mathcal O}\), as is seen on the
simple factors. Equivalently the signed length comparison is
\[
 r_{\mathcal O}m_{\ell,e}
       \le\sum_{\psi(c)=(-1)^{e+1}} V((\psi\kappa^e)^{-1})
\]
in extended lengths; the character bound makes the right side finite.
For any nonzero \(a\in\mathcal O\), the normalized absolute value
satisfies \(v(a)\log\ell=-r_{\mathcal O}\log|a|_\ell\).
Multiplying the displayed inequality by \(\log\ell\), applying
\eqref{eq:characterwise-bern}, and cancelling \(r_{\mathcal O}>0\)
give the inequality used before
\eqref{eq:galois-good-prime-sum}. Distinct characters of \(G\) pull
back to distinct Dirichlet characters because the cyclotomic action
map onto \(G\) is surjective.

Sections~\ref{sec:minus-derivatives} and
\ref{sec:minus-cyclotomic-elements} prove the cyclotomic
construction and the two local derivative identities.
The simultaneous Chebotarev choice is proved in
Section~\ref{sec:minus-frobenius}. Together with the two duality
theorems and the preceding descent these yield
\eqref{eq:characterwise-bern}.
\footnote{\href{https://github.com/Deicyde/riemann-hypothesis/blob/40e56bab026552e089d7afcfeff08d2defb9c005/src/Development/AutomorphicArithmetic/AutomorphicVarieties/PrimitiveFrobeniusFixed/Iv.lean}{Precise input to the characterwise theorem},
\href{https://github.com/Deicyde/riemann-hypothesis/blob/40e56bab026552e089d7afcfeff08d2defb9c005/src/Development/AutomorphicArithmetic/AutomorphicVarieties/PrimitiveFrobeniusFixed/OverIff.lean}{global and local reciprocity},
\href{https://github.com/Deicyde/riemann-hypothesis/blob/40e56bab026552e089d7afcfeff08d2defb9c005/src/Development/AutomorphicArithmetic/AutomorphicVarieties/PrimitiveFrobeniusFixed/ZeNorm.lean}{finite Bernoulli value and Euler factor},
\href{https://github.com/Deicyde/riemann-hypothesis/blob/40e56bab026552e089d7afcfeff08d2defb9c005/src/Development/AutomorphicArithmetic/AutomorphicVarieties/PrimitiveFrobeniusFixed/RatPlace.lean}{the exceptional case},
\href{https://github.com/Deicyde/riemann-hypothesis/blob/40e56bab026552e089d7afcfeff08d2defb9c005/src/Development/AutomorphicArithmetic/AutomorphicVarieties/PrimitiveFrobeniusFixed/NuSysEx.lean}{cyclotomic derivatives},
\href{https://github.com/Deicyde/riemann-hypothesis/blob/40e56bab026552e089d7afcfeff08d2defb9c005/src/Development/AutomorphicArithmetic/AutomorphicVarieties/PrimitiveFrobeniusFixed/AddPowNsmul.lean}{Frobenius selection},
\href{https://github.com/Deicyde/riemann-hypothesis/blob/40e56bab026552e089d7afcfeff08d2defb9c005/src/Development/AutomorphicArithmetic/AutomorphicVarieties/PrimitiveFrobeniusFixed/ChiDvd.lean}{extension to the dual cocycle},
\href{https://github.com/Deicyde/riemann-hypothesis/blob/40e56bab026552e089d7afcfeff08d2defb9c005/src/Development/AutomorphicArithmetic/AutomorphicVarieties/PrimitiveFrobeniusFixed/FrobFix.lean}{length induction},
\href{https://github.com/Deicyde/riemann-hypothesis/blob/40e56bab026552e089d7afcfeff08d2defb9c005/src/Development/AutomorphicArithmetic/AutomorphicVarieties/PrimitiveFrobeniusFixed/PrimCase.lean}{primitive conductor comparisons},
\href{https://github.com/Deicyde/riemann-hypothesis/blob/40e56bab026552e089d7afcfeff08d2defb9c005/src/Development/AutomorphicArithmetic/AutomorphicVarieties/OddResidueBounds/AtPrime.lean}{the signed character sum}.}
